\clearpage\fancyhead[L]{\small GAUSS TO BODE.}\hypertarget{native-362-lower}{}\label{native-362-lower}
\begin{center}\textbf{GAUSS TO BODE.}\par\vspace{4pt}\hrule\vspace{4pt}\fontsize{10}{12}\selectfont Astronomisches Jahrbuch für 1815. pp. 245\;..\;249. Berlin 1812.\par\vspace{4pt}\hrule\vspace{4pt}\end{center}

\begin{flushright}Göttingen, 22 August 1812.\end{flushright}
\par\indent Unless I am mistaken, I have already told you about the first sighting of Pallas. There have been only a few observations here this time. Around the time of opposition the
\SourceLeaf{363}planet was too high in the sky at our site to observe with the usual instrument and ring micrometer. I therefore had to use another, inferior instrument. Even then, poor weather and moonlight allowed me only one observation. Here are all my observations:

\begin{center}\fontsize{10}{12}\selectfont\setlength{\tabcolsep}{4pt}\renewcommand{\arraystretch}{1.15}\begin{tabular}{l|c|c|c}
& Mean time & Apparent R.A. & Northern decl.\\\hline
April 9 & $11^{\mathrm u}\;32^{\mathrm m}\;27^{\mathrm s}$ & $269^\circ\;5'\;58''4$ & $16^\circ\;38'$\\
May 2 & $10\;36\;6$ & $268\;42\;5.3$ & $21\;4\;31''7$\\
— 4 & $10\;44\;44$ & $268\;31\;38.0$ & $21\;25\;4.3$\\
June 7 & $11\;13\;51$ & $262\;50\;51.1$ & $25\;2\;56.6$\\
\end{tabular}\end{center}

\par\indent Apart from two earlier observations by Dr. \textsc{Olbers}, the only observations from elsewhere that have been sent to me are four meridian observations by Dr. \textsc{Burckhardt} at the observatory of the military school in Paris. I had them printed in no. 127 of our Gelehrte Anzeigen [p. 356 of this volume]. Mr. \textsc{Nicolai}, whose skill as a calculator you already know, calculated the opposition from these data:

\Needspace{10\baselineskip}1812 June 10. $3^{\mathrm u}\;2^{\mathrm m}\;48^{\mathrm s}$ Mean time in Göttingen

\begin{center}\fontsize{10}{12}\selectfont\setlength{\tabcolsep}{2pt}\renewcommand{\arraystretch}{1.15}\begin{tabular}{lcr}
True longitude & $.\;.\;.\;.\;.\;.\;.\;.\;.\;.\;.$ & $259^\circ\;27'\;51''7$\\
True geoc. latitude & $.\;.\;.\;.\;.\;.\;.\;.\;.$ & $48\;16\;26.1\;\mathrm{N}.$\\
\end{tabular}\end{center}

\Needspace{13\baselineskip}\par\indent With only the following small corrections made to the latest elements,

\begin{center}\fontsize{10}{12}\selectfont\setlength{\tabcolsep}{4pt}\renewcommand{\arraystretch}{1.15}\begin{tabular}{rl}
$-13''29$ & to the epoch of longitude 1810\\
$-0.01200$ & to the daily motion\\
$-17.57$ & to the longitude of perihelion\\
$-4.09$ & to the angle of eccentricity\\
$+0''80$ & to the longitude of the node\\
$-1.21$ & to the inclination of the orbit\\
$+0.0000045$ & to the logarithm of the semimajor axis;\\
\end{tabular}\end{center}

\Needspace{17\baselineskip}\noindent the closest possible agreement was as follows.

\begin{center}\fontsize{10}{12}\selectfont\setlength{\tabcolsep}{6pt}\renewcommand{\arraystretch}{1.15}\begin{tabular}{c|c|c}
& \multicolumn{2}{c}{Difference}\\
Oppos. & Mean longitude & Helioc. latitude\\\hline
$1803$ & $-5''06$ & $-2''78$\\
$1804$ & $-12.95$ & $-9.66$\\
$1805$ & $-15.20$ & $-0.67$\\
$1807$ & $+8.28$ & $-5.77$\\
$1808$ & $+16.04$ & $+11.67$\\
$1809$ & $+13.00$ & $-1.48$\\
$1811$ & $+11.73$ & $-9.53$\\
$1812$ & $-25.92$ & $+6.15$\\
\end{tabular}\end{center}

\par\indent I attribute these differences chiefly to the perturbations due to Saturn and Mars, since only the perturbations due to Jupiter have been taken into account so far. The latter, however, are by far the largest. The first calculation of them gave me about 400 equations. In a second, more complete calculation on which I have made a small start, about 1000 will be developed that exceed $0''1$. According to Mr. \textsc{Nicolai}’s calculation, the next opposition will occur

\Needspace{10\baselineskip}1813 August 18. $8^{\mathrm u}\;50^{\mathrm m}\;5^{\mathrm s}$ Mean time in Göttingen

\begin{center}\fontsize{10}{12}\selectfont\setlength{\tabcolsep}{2pt}\renewcommand{\arraystretch}{1.15}\begin{tabular}{lcr}
True longitude & $.\;.\;.\;.\;.\;.\;.\;.\;.\;.\;.\;.$ & $325^\circ\;24'\;18''$\\
True geocentric latitude & $.\;.\;.\;.\;.\;.\;.$ & $24\;37\;24\;\mathrm{N}..$\\
\end{tabular}\end{center}

\Needspace{17\baselineskip}\SourceLeaf{364}\par\indent Observations of \emph{Juno} were extremely difficult this year. Even finding it in the Milky Way, which teems with faint stars, took a great deal of effort because the planet was so faint. It was only about the 11th magnitude, and even later did not reach the 10th. Once I had found it, however, I was fortunate enough to make 4 good observations, which I give you here.

\begin{center}\fontsize{10}{12}\selectfont\setlength{\tabcolsep}{4pt}\renewcommand{\arraystretch}{1.15}\begin{tabular}{l|c|c|c}
1812 & Mean time & Apparent R.A. & Declination\\\hline
Jan. 7 & $12^{\mathrm u}\;21^{\mathrm m}\;0^{\mathrm s}$ & $282^\circ\;21'\;24''5$ & $4^\circ\;41'\;44''6\;\mathrm S.$\\
13 & $12\;51\;12$ & $281\;16\;32.9$ & $4\;33\;39.6$\\
July 12 & $11\;35\;28$ & $275\;5\;10.3$ & $5\;6\;11.2$\\
13 & $11\;28\;12$ & $274\;52\;49.9$ & $5\;9\;20.0$\\
\end{tabular}\end{center}

\Needspace{15\baselineskip}\par\indent Mr. \textsc{Wachter} determined the opposition from these observations by a method of my own, which allows observations farther away in time than the usual procedure does. I must save its explanation for another occasion. Here is his result:

1812 June 29. $19^{\mathrm u}\;51^{\mathrm m}\;43^{\mathrm s}$ Mean time in Göttingen

\begin{center}\fontsize{10}{12}\selectfont\setlength{\tabcolsep}{2pt}\renewcommand{\arraystretch}{1.15}\begin{tabular}{lcr}
True longitude & $.\;.\;.\;.\;.\;.\;.\;.\;.\;.\;.\;.$ & $278^\circ\;15'\;21''3$\\
True geocentric latitude & $.\;.\;.\;.\;.\;.\;.$ & $18\;36\;56.9$\\
\end{tabular}\end{center}

\Needspace{14\baselineskip}\par\indent The differences between my latest elements and the observations were determined in the usual way. Mr. \textsc{Wachter} calculated the following:

\begin{center}\fontsize{10}{12}\selectfont\setlength{\tabcolsep}{7pt}\renewcommand{\arraystretch}{1.15}\begin{tabular}{l|c|c}
& R.A. & Decl.\\\hline
June 7 & $+21''8$ & $-35''9$\\
13 & $+31.0$ & $-34.7$\\
July 12 & $+22.4$ & $-46.9$\\
13 & $+26.2$ & $-45\;0$\\
\end{tabular}\end{center}

\Needspace{16\baselineskip}\par\indent The same skilled calculator also fitted the elements of the 4 most recent oppositions by the method I explained in volume I of the Neuen Comment. of our Royal Society [p. 1 of this volume], and found the following results:

\begin{center}\fontsize{9}{11}\selectfont\setlength{\tabcolsep}{2pt}\renewcommand{\arraystretch}{1.15}\begin{tabular}{lcr}
Epoch, Göttingen meridian & $.\;.\;.\;.\;.\;.\;.\;.$ & $177^\circ\;48'\;21''0$\\
Daily mean tropical motion & $.\;.\;.\;.\;.$ & $813''25748$\\
Longitude of perihelion 1811 & $.\;.\;.\;.\;.\;.\;.\;.\;.\;.\;.$ & $53^\circ\;15'\;10''1$\\
Longitude of $\AscendingNode$ & $.\;.\;.\;.\;.\;.\;.\;.\;.\;.\;.\;.\;.\;.$ & $171\;9\;16.7$\\
Inclination of the orbit & $.\;.\;.\;.\;.\;.\;.\;.\;.\;.\;.\;.\;.$ & $13\;4\;17.2$\\
Logarithm of the semimajor axis & $.\;.\;.\;.\;.\;.$ & $0.4265679$\\
\end{tabular}\end{center}

\par\indent Finally, he also calculated the ephemeris for Juno for next year, and I have the pleasure of enclosing a copy. It does not appear that Juno has been observed anywhere else this year.

\par\indent Professor \textsc{Harding} observed Vesta a few times with the ring micrometer. The ephemeris gave the right ascension about $11'$ too small. No doubt its $\text{\AstroSymbols ☍}$ can be determined very well from meridian observations. I have chosen to have no new calculations made before then. It will also be enough for a new ephemeris to begin around the end of 1813, so none will yet be needed for the yearbook for 1815.

\SourceLeaf{365}\par\indent Ceres was also observed here a few times on and off the meridian around the time of $\text{\AstroSymbols ☍}$. No calculations have yet been made from these observations, however.

\par\indent Professor \textsc{Harding} and I spent an evening looking in vain for the new comet announced in the newspapers.


